\documentclass{exam}

\usepackage[margin=1in]{geometry}

\usepackage{comment}
\usepackage{graphicx}
\usepackage{listings}

\pagestyle{headandfoot}
\firstpageheader{CSCI 221}{}{C++ Intro Assignment}
\runningheader{CSCI 221}{}{C++ Intro Assignment}


\begin{document}
\title{C++ Intro Assignment (Evaluation)}
\author{CSCI 221: Computer Science Fundamentals 2}
\date{Spring 2026}
\maketitle


This assignment is an opportunity to test your understanding of C++. 
Feel free to collaborate with others and use resources, but all code and writeups must be your own. 
To submit your assignment, please submit your code, makefile, and writeup. 
For full credit, your code should check for invalid input. 

\textbf{Due Date:}
Friday, April 24th at 1:00 pm. 

\begin{questions}

\question[15]
\textbf{Cars.}
\begin{parts}
\part[5]
Write the code for a Car class. 
It should represent car objects that move from location to location using fuel from their gas tank. 
Cars should have variables representing their position in a two-dimensional space, their fuel efficiency in miles per gallon, the size of their fuel tank, and how much fuel they currently have. 
Cars should also have a method \texttt{move\_to} that tries to move to new coordinates. 
If the car has enough fuel to drive to the new coordinates in a straight line, it does so and returns true. 
If not, the car does not move and returns false. 

\part[3]
Write a function that given an array of Cars and a point in two-dimensional space, returns an array containing a copy of each of the cars that can move to that space. 
The input array and the Cars in it should not be modified. 
The Cars in the output array should have the \texttt{move\_to} method called on them prior to returning. 
Make sure you are also returning the length of the output array. 

\part[2]
Write the code for a GasStation class. 
It should contain a location and a price in dollars per gallon for refueling cars.

\part[5]
Write a modified version of your function for moving an array of cars that also takes an array of GasStations. 
Cars that cannot directly move to the target location should check to see if they can move to the target point using the GasStations. 
If this is possible, they should be added to the output list as if they move there using that path. 
For simplicity, assume that the array of GasStations is sorted in ascending order of proximity to the target point and that traveling to a GasStation that is further away cannot help a Car reach the target point. 
In addition to the list of Cars being output, you should output a list of the costs for each car to fill up its tank on the route to the location. 
(This is zero for the cars that do not need to go to a GasStation.)

\end{parts}

\question[25]
\textbf{SortedLists.}
Write the code for a SortedList class. 
This class should contain a linked list of strings in ascending sorted order. 
There are two possible ways the list can be sorted:
\begin{enumerate}
\item By string length. 
\item By individual characters. 
Starting with the first, one character from each string is compared using the numeric value of the ASCII encoding. 
If the strings have the same character, compare them on the next. 
\end{enumerate}
Ties can be broken in any order. 

Your class should have the following functionality. 
\begin{parts}
\part[1]
The ability to create new instances of the class that represent empty lists. 

\part[1]
The ability to create (deep) copies of a SortedList. 

\part[1]
The ability to create new instances of the class that specify which sorting method to use. 

\part[3]
A method to add elements to the list in sorted order. 

\part[3]
A method to remove all instances of an element from the list. 

\part[1]
A method to remove all values from the list. 

\part[1]
A method that returns whether or not the list is currently sorted by string length. 

\part[1]
A method that returns whether or not the list is currently sorted by individual characters. 

\part[1]
A method that returns the length of the list. 

\part[2]
A method that returns an array containing the values in the list in current sorted order. 

\part[2]
A method that returns the value in the $i$th position in current sorted order. 

\part[3]
A method that changes the list to be sorted by string length. 

\part[3]
A method that changes the list to be sorted by individual characters. 

\part[1]
The class should not cause any memory leaks. 

\part[1]
The class should have functions properly marked as constant. 
\end{parts}

\question[10]
\textbf{Writeup.}
In addition to your code, you should submit a brief (~1-2 pages) writeup that describes the state of your code. 
You should discuss what is working, what is not, and any known errors. 
In this discussion, you should reference the tests you did and how they support your claims. 
You will need to provide a brief description of your tests for you to reference. 
Your writeup should be in a separate .txt or .pdf file. 

\end{questions}

\end{document}